\chapter{Market Quality and Regulation}\label{ch:mx:market-quality-and-regulation}

Twenty years of automation have narrowed spreads and multiplied messages, and regulators still disagree about whether markets are better for it.
The evidence is abundant; what it measures, and what it cannot, is the subject of this chapter. It defines the measures, reads the main studies of
algorithmic and \omterm{def:mx:market-quality-and-regulation:hft}{high-frequency trading}, the flash crash and latency races, and then evaluates a planted rule on a simulated panel of stock-days, where
the rule's true effect is known, to show what a \omterm{def:mx:tick-size-queues-and-priority:did}{difference-in-differences} estimate recovers and what a before-and-after comparison gets wrong.

\section{Measuring liquidity and market quality}

\begin{definition}[Market quality]\label{def:mx:market-quality-and-regulation:quality}
\emph{Market quality}\index{market quality} is the set of properties by which a market serves those who trade in it: the cost of trading (quoted,
effective and realised spreads, price impact), the quantity available (depth, resilience), and the informativeness of prices (how close they stay to
a random walk, how fast they incorporate information).
\end{definition}

The measures are those of earlier chapters and books: the \omterm{def:mx:decomposing-the-spread:quoted}{quoted spread} (chapter~5), the effective and realised spreads and their difference, the
price impact (One Quant Book 1, chapter~10), depth and resilience (chapter~3), Amihud's illiquidity ratio (One Quant Book 7, chapter~7), and the variance
ratio, the variance of returns over a long interval against the sum of the variances over the short intervals within it (one for a random walk, below one
when prices bounce between bid and ask). \texttt{firm\_mktquality.day\_measures} computes them for a stock-day from a top-of-book path, the trades and
the number of order messages.

\begin{definition}[Order-to-trade ratio]\label{def:mx:market-quality-and-regulation:otr}
The \emph{order-to-trade ratio}\index{order-to-trade ratio} of a participant or a stock is the number of order messages (new orders, modifications,
cancellations) per trade over a period.
\end{definition}

\section{Algorithmic and high-frequency trading in the evidence}

\begin{definition}[Algorithmic trading, high-frequency trading]\label{def:mx:market-quality-and-regulation:hft}
\emph{Algorithmic trading}\index{algorithmic trading} is trading in which a computer algorithm determines the parameters of orders (whether to send
them, their timing, price or quantity, how to manage them afterwards) with little or no human intervention. \emph{High-frequency
trading}\index{high-frequency trading} is algorithmic trading that relies on infrastructure minimising latency (co-location, proximity hosting or
high-speed direct access), decides without human intervention for individual orders, and sends many orders, quotes or cancellations within the day.
\end{definition}

The definitions follow MiFID~II's article~4(1)(39) and (40), which exclude systems that only route orders or process them without deciding their
parameters. The evidence is mostly favourable to liquidity in normal times. Hendershott, Jones and Menkveld (2011) used the New York Stock Exchange's
automation of quote dissemination in 2003, which increased \omterm{def:mx:market-quality-and-regulation:hft}{algorithmic trading}, as an instrument: for large stocks, \omterm{def:mx:market-quality-and-regulation:hft}{algorithmic trading} narrowed
spreads, reduced adverse selection and made quotes more informative. Brogaard, Hendershott and Riordan (2014) found that high-frequency traders
facilitate price efficiency by trading in the direction of permanent price changes and against transitory pricing errors, mainly through their
liquidity-demanding orders, while their liquidity-supplying orders are adversely selected. Menkveld (2013) characterised one large high-frequency
trader: it lost on its inventory, earned the spread, took part in 8.1\% of trades on the incumbent market and 64.4\% on an entrant, and was passive in
four trades out of five.

\section{The flash-crash literature}

\begin{definition}[Hot-potato trading]\label{def:mx:market-quality-and-regulation:hotpotato}
\emph{Hot-potato trading}\index{hot-potato trading} is intermediaries passing the same positions rapidly back and forth among themselves, which
inflates volume without transferring risk to anyone who wants to hold it.
\end{definition}

On 6 May 2010 (One Quant Book 1, chapter~31), a mutual fund complex began to sell 75\,000 E-mini contracts, about USD~4.1 billion, with an algorithm set to
trade 9\% of the previous minute's volume without regard to price or time: chapter~16's \omterm{def:mx:benchmark-algorithms:vwap}{participation algorithm}, with the feedback that chapter
described. The CFTC and SEC staff report found that, lacking fundamental buyers, high-frequency traders began to buy and resell contracts to each
other; between 2:45:13 and 2:45:27 they traded over 27\,000 contracts, about 49\% of volume, while buying only about 200 net: a hot-potato volume effect
that the volume-following algorithm read as liquidity. Kirilenko, Kyle, Samadi and Tuzun (2017), with the audit trail of the E-mini, found that the
most active non-designated intermediaries did not change their trading pattern when prices fell. Volume is not liquidity: a market-quality measure that
counts trades would have scored the flash crash's worst minutes as its most liquid.

\section{Latency races and speed bumps in the evidence}

Aquilina, Budish and O'Neill (2022) used exchange message data, which record failed attempts to trade or cancel as well as successful ones, to see
latency-arbitrage races directly: for FTSE~100 stocks, about one race per minute per symbol, a modal race lasting 5 to 10 millionths of a second, and
about a fifth of all trading volume. Chapter~10 showed that \omterm{def:mx:auction-mechanics-and-theory:fba}{frequent batch auctions} remove the race by design; speed bumps and the economics of the race
itself are One Quant Book 11's subject. For this chapter the lesson is methodological: the races were invisible in the order-book data most studies
use, and a measure can only see what its data record.

\section{Rules on algorithms and order-to-trade ratios}

The SEC's 2010 concept release on equity market structure asked for comment on \omterm{def:mx:market-quality-and-regulation:hft}{high-frequency trading}, order routing, market data linkages and
undisplayed liquidity; MiFID~II followed with its definitions and with a requirement on venues to track each member's ratio of unexecuted orders to
transactions.

\begin{dated}{2026-09}{Order-to-trade ratios in the EU}\label{dat:mx:market-quality-and-regulation:otr}
Commission Delegated Regulation (EU) 2017/566 (RTS~9 under MiFID~II) requires trading venues to calculate the ratio of unexecuted orders to transactions
of each member and participant, for every financial instrument traded in a continuous order book, quote-driven or hybrid system; it applies from the
date MiFID~II applied. The text quoted is the regulation as adopted.
\end{dated}

Whether such a rule improves the market is an empirical question, and the simulation answers it where the truth is known. Twenty stocks of
different activity trade for sixteen ten-minute days in \texttt{firm.agentmkt}. From day~8 two things happen: every stock's market becomes more
volatile and its liquidity providers send a quarter fewer orders (a market-wide change, the confounder); and every other stock, ordered by activity,
comes under an order-to-trade cap that makes its liquidity providers send 30\% fewer orders and cancel less. Each treated stock-day after the event is
also run without the cap, with the same order flow, so the rule's effect is measured directly.

\begin{omfigure}
\begin{tikzpicture}
\begin{axis}[width=9.6cm,height=5.4cm,xlabel={\small day relative to the rule},ylabel={\small quoted spread (bp)},
  tick label style={font=\footnotesize},grid=major,grid style={gray!25},xmin=-8.5,xmax=7.5,ymin=1,ymax=3.6,
  legend columns=2,legend style={font=\footnotesize,at={(0.5,-0.3)},anchor=north,draw=none,fill=none,
  /tikz/every even column/.append style={column sep=8pt}},table/col sep=comma]
\addplot[omProp,very thick,mark=*,mark size=1.4pt] table[x=day,y=treated]{figdata/microstructure/24-market-quality-and-regulation/daily.csv};
\addplot[omDef,very thick,mark=square*,mark size=1.4pt] table[x=day,y=control]{figdata/microstructure/24-market-quality-and-regulation/daily.csv};
\addplot[gray,dashed] coordinates {(-0.5,1) (-0.5,3.6)};
\legend{treated stocks,control stocks}
\end{axis}
\end{tikzpicture}

\omcaption{The mean \omterm{def:mx:decomposing-the-spread:quoted}{quoted spread} of the ten treated and ten control stocks, day by day. Both rise at the event (the market-wide change); the treated
rise more (the rule). Data: \texttt{mx\_mq.study}.}\label{fig:mx:market-quality-and-regulation:daily}
\end{omfigure}

The treated stocks' \omterm{def:mx:decomposing-the-spread:quoted}{quoted spread} rose by 1.60 basis points (standard error 0.08) from before to after; the \omterm{def:mx:tick-size-queues-and-priority:did}{difference-in-differences} estimate
(chapter~6's method, with stock and day effects and errors clustered by stock, \texttt{firm\_mktquality.did}) is 0.89 (0.09); the
truth is 0.87 (0.03). The before-and-after comparison attributes the market-wide change of 0.71 basis points to the rule
(\cref{tab:mx:market-quality-and-regulation:did}).

\omcode{code/firm/mktquality/firm_mktquality.py}{61}{71}{The difference-in-differences estimate with two-way fixed effects: the treated-after indicator, stock and day dummies, and standard errors clustered by stock.}\label{lst:mx:market-quality-and-regulation:did}

\begin{table}[ht]
\centering
\small
\begin{tabular}{@{}lrrr@{}}
\toprule
measure & before-after (s.e.) & \omterm{def:mx:tick-size-queues-and-priority:did}{difference-in-differences} (s.e.) & truth (s.e.) \\
\midrule
\omterm{def:mx:decomposing-the-spread:quoted}{quoted spread} (bp) & 1.60 (0.08) & 0.89 (0.09) & 0.87 (0.03) \\
effective spread (bp) & 2.74 (0.12) & 1.31 (0.14) & 1.36 (0.06) \\
realised spread (bp) & 1.81 (0.08) & 1.06 (0.10) & 1.15 (0.07) \\
price impact (bp) & 0.93 (0.09) & 0.25 (0.10) & 0.21 (0.07) \\
depth (shares) & $-138$ (8) & $-20$ (13) & $-31$ (2) \\
variance ratio & $-0.19$ (0.02) & $-0.13$ (0.06) & $-0.14$ (0.03) \\
\omterm{def:mx:market-quality-and-regulation:otr}{order-to-trade ratio} & $-2.64$ (0.05) & $-1.18$ (0.07) & $-1.19$ (0.02) \\
\bottomrule
\end{tabular}
\caption{The rule's effect on the treated stocks: the before-and-after change, the \omterm{def:mx:tick-size-queues-and-priority:did}{difference-in-differences} estimate and the true effect from the
same days run without the rule. Data: \texttt{mx\_mq.study}.}\label{tab:mx:market-quality-and-regulation:did}
\end{table}

The rule did what it was meant to do (1.2 fewer messages per trade) and made the market worse on every other measure: wider spreads, less depth, more
bid-ask bounce. The event study (\cref{fig:mx:market-quality-and-regulation:event}) checks the design: before the event the treated and control stocks'
spreads move together (the seven coefficients are within 0.07 basis points of zero), and after it the gap opens at once and stays.

\begin{omfigure}
\begin{tikzpicture}
\begin{axis}[width=9.6cm,height=5.4cm,xlabel={\small day relative to the rule},ylabel={\small treated minus control (bp)},
  tick label style={font=\footnotesize},grid=major,grid style={gray!25},xmin=-8.5,xmax=7.5,ymin=-0.4,ymax=1.6,table/col sep=comma]
\addplot[omDef,very thick,mark=*,mark size=1.4pt,error bars/.cd,y dir=both,y explicit] table[x=day,y=coef,y error expr=1.96*\thisrow{se}]{figdata/microstructure/24-market-quality-and-regulation/event.csv};
\addplot[black,dashed,domain=-8.5:7.5] {0};
\end{axis}
\end{tikzpicture}

\omcaption{Event study of the \omterm{def:mx:decomposing-the-spread:quoted}{quoted spread}: the treated-minus-control difference each day relative to the day before the rule, with stock and day
effects and 95\% intervals clustered by stock. Data: \texttt{mx\_mq.study}.}\label{fig:mx:market-quality-and-regulation:event}
\end{omfigure}

\section{Tutorial: did the rule make the market better?}\label{tut:mx:market-quality-and-regulation:tut}

\begin{tutorial}
\textbf{Goal.} Build a stock-day panel of market-quality measures in the simulated market, plant a rule on half the stocks, and estimate it by
\omterm{def:mx:tick-size-queues-and-priority:did}{difference-in-differences} against the before-and-after comparison and the truth.
\textbf{End state:} \cref{fig:mx:market-quality-and-regulation:daily,fig:mx:market-quality-and-regulation:event},
\cref{tab:mx:market-quality-and-regulation:did}.
\begin{tutsteps}
\item \textbf{Measures.} \texttt{firm\_mktquality.day\_measures(top, trades, messages)}.
\item \textbf{Panel.} \texttt{mx\_mq.config(i, d, rule)}, \texttt{session\_day}, \texttt{panel()}.
\item \textbf{Estimate.} \texttt{did}, \texttt{before\_after}, \texttt{event\_study}; \texttt{study()}; draw with \texttt{fig\_mq.py}.
\end{tutsteps}
\textbf{What to change next.} Treat the ten least active stocks instead of every other one and watch the parallel-trends assumption fail
(exercise~7); compute the measures on SEC MIDAS data for a real rule change.
\end{tutorial}

\section{Build: market quality}\label{bld:mx:market-quality-and-regulation:mktquality}

\begin{build}
\bfield{Purpose} Measuring markets and rules: the panel measures and estimators behind this chapter and chapter~25's circuit breakers.
\bfield{Interface} \texttt{day\_measures(top, trades, messages, h, grid, px\_per\_ccy)}, \texttt{did(y, treated, post, unit, period)},
\texttt{before\_after(y, treated, post, unit)}, \texttt{event\_study(y, treated, period, unit, base)}.
\bfield{Rules} Spreads in basis points of the mid; firm.tape's trade and top-of-book fields; errors clustered by the unit of treatment.
\bfield{Acceptance tests} \texttt{code/firm/mktquality/tests/}: the measures on a hand market; a planted effect recovered by \omterm{def:mx:tick-size-queues-and-priority:did}{difference-in-differences}
and missed by before-and-after; a flat pre-period in the event study.
\bfield{Stretch} Resilience; synthetic controls; staggered adoption.
\end{build}

\begin{omsources}
\item CFTC and SEC staff, ``Findings regarding the market events of May 6, 2010'', 2010.
\item SEC, Concept release on equity market structure, Release No.~34-61358, 2010.
\item T.~Hendershott, C.~M.~Jones and A.~J.~Menkveld, ``Does algorithmic trading improve liquidity?'', \emph{Journal of Finance} 66(1), 2011.
\item A.~J.~Menkveld, ``High frequency trading and the new market makers'', \emph{Journal of Financial Markets} 16(4), 2013.
\item J.~Brogaard, T.~Hendershott and R.~Riordan, ``High-frequency trading and price discovery'', \emph{Review of Financial Studies} 27(8), 2014.
\item A.~Kirilenko, A.~S.~Kyle, M.~Samadi and T.~Tuzun, ``The flash crash: high-frequency trading in an electronic market'', \emph{Journal of
Finance} 72(3), 2017.
\item M.~Aquilina, E.~Budish and P.~O'Neill, ``Quantifying the high-frequency trading `arms race''', \emph{Quarterly Journal of Economics} 137(1),
2022.
\item Directive 2014/65/EU (MiFID~II), article 4(1)(39)--(40); Commission Delegated Regulation (EU) 2017/566.
\end{omsources}

\section{Exercises}

\begin{exercise}[$\star$]\label{exo:mx:market-quality-and-regulation:1}
From the four group means of the \omterm{def:mx:decomposing-the-spread:quoted}{quoted spread} (treated 1.58 before and 3.18 after, control 1.56 and 2.27), compute the \omterm{def:mx:tick-size-queues-and-priority:did}{difference-in-differences}.
\end{exercise}

\begin{exercise}[$\star$]\label{exo:mx:market-quality-and-regulation:2}
A stock-day has 900 order messages and 250 trades. What is its \omterm{def:mx:market-quality-and-regulation:otr}{order-to-trade ratio}, and what would the rule's effect of $-1.18$ make it?
\end{exercise}

\begin{exercise}[$\star$]\label{exo:mx:market-quality-and-regulation:3}
Why is a variance ratio below one a sign of bid-ask bounce, and what did the rule do to it?
\end{exercise}

\begin{exercise}[$\star\star$]\label{exo:mx:market-quality-and-regulation:4}
Why do the standard errors cluster by stock and not by stock-day?
\end{exercise}

\begin{exercise}[$\star\star$]\label{exo:mx:market-quality-and-regulation:5}
Why did high-frequency traders' volume during the flash crash mislead a \omterm{def:mx:benchmark-algorithms:vwap}{participation algorithm}?
\end{exercise}

\begin{exercise}[$\star\star$]\label{exo:mx:market-quality-and-regulation:6}
Why can order-book data not show latency races, and what data can?
\end{exercise}

\begin{exercise}[$\star\star\star$]\label{exo:mx:market-quality-and-regulation:7}
\emph{Coding.} Treat the ten least active stocks instead of every other one. What does \omterm{def:mx:tick-size-queues-and-priority:did}{difference-in-differences} estimate for the \omterm{def:mx:decomposing-the-spread:quoted}{quoted spread}, and
why does it miss the truth?
\end{exercise}

\begin{exercise}[$\star\star\star$]\label{exo:mx:market-quality-and-regulation:8}
\emph{Find the flaw.} ``Since the rule, spreads on the stocks it covers have risen by 1.6 basis points: the rule cost investors 1.6 basis points.''
\end{exercise}

\section{Problem: Did the Rule Make the Market Better?}

\begin{problem}[{Weekend problem --- did the rule make the market better?}]\label{pb:mx:market-quality-and-regulation:1}
A regulator capped \omterm{def:mx:market-quality-and-regulation:otr}{order-to-trade ratios} on half the stocks of a market. You are asked to evaluate it.

\medskip\noindent\textbf{Part I --- Measures.}
\begin{enumerate}
\item Define \omterm{def:mx:market-quality-and-regulation:quality}{market quality} and list its measures.
\item How are the effective and realised spreads and the price impact related?
\item Define the \omterm{def:mx:market-quality-and-regulation:otr}{order-to-trade ratio}, \omterm{def:mx:market-quality-and-regulation:hft}{algorithmic trading} and \omterm{def:mx:market-quality-and-regulation:hft}{high-frequency trading}.
\item What does the variance ratio measure?
\end{enumerate}

\medskip\noindent\textbf{Part II --- The evidence.}
\begin{enumerate}[resume]
\item What did Hendershott, Jones and Menkveld find, and how did they identify it?
\item What did Brogaard, Hendershott and Riordan find about \omterm{def:mx:fragmentation-and-routing:discovery}{price discovery}?
\item Describe the flash crash's sell programme and the hot-potato volume.
\item What did Aquilina, Budish and O'Neill measure, and why could they?
\end{enumerate}

\medskip\noindent\textbf{Part III --- The simulated rule.}
\begin{enumerate}[resume]
\item Describe the panel, the confounder and the rule.
\item How is the truth measured?
\item Write the two-way fixed-effects regression.
\item \emph{State the named result}: the \omterm{def:mx:tick-size-queues-and-priority:did}{difference-in-differences} estimate of the rule's effect on \omterm{def:mx:decomposing-the-spread:quoted}{quoted spreads} with its standard error, against the
before-after estimate that ignores the control group.
\item What does the event study show before the event, and why does it matter?
\end{enumerate}

\medskip\noindent\textbf{Part IV --- Judgement.}
\begin{enumerate}[resume]
\item Did the rule make the market better? Use the table.
\item What happens if the treated stocks are the least active ones?
\item Which measure moved most in relative terms, and why?
\item What would you need to evaluate such a rule on real data?
\item Why is volume a poor measure of liquidity?
\item What does the dated box say about \omterm{def:mx:market-quality-and-regulation:otr}{order-to-trade ratios} in the EU?
\item In one sentence: what does a control group buy an evaluation?
\end{enumerate}
\end{problem}

\section{Interview questions}

\begin{interviewq}[$\star$ \iqroles{trader}]\label{iq:mx:market-quality-and-regulation:1}
Spreads in your stock doubled after a rule change. How do you know the rule did it?
\end{interviewq}

\begin{interviewq}[$\star\star$ \iqroles{researcher}]\label{iq:mx:market-quality-and-regulation:2}
Explain \omterm{def:mx:tick-size-queues-and-priority:did}{difference-in-differences} and its key assumption. How would you test it?
\end{interviewq}

\begin{interviewq}[$\star\star$ \iqroles{researcher}]\label{iq:mx:market-quality-and-regulation:3}
Does \omterm{def:mx:market-quality-and-regulation:hft}{high-frequency trading} improve \omterm{def:mx:market-quality-and-regulation:quality}{market quality}? What does the evidence say, and what does it not?
\end{interviewq}

\begin{interviewq}[$\star\star$ \iqroles{risk}]\label{iq:mx:market-quality-and-regulation:4}
What lessons does the flash crash hold for an \omterm{def:mx:benchmark-algorithms:algo}{execution algorithm}?
\end{interviewq}

\begin{interviewq}[$\star\star$ \iqroles{developer}]\label{iq:mx:market-quality-and-regulation:5}
Your firm must stay under a venue's \omterm{def:mx:market-quality-and-regulation:otr}{order-to-trade ratio}. How would you monitor and enforce it in your trading systems?
\end{interviewq}

\begin{interviewq}[$\star\star\star$ \iqroles{researcher}]\label{iq:mx:market-quality-and-regulation:6}
How would you measure latency races on a venue, and what would you need from it?
\end{interviewq}
